% house_solmacros.tex: HOUSE COPY of de-docs/common/de_style.tex minus \input{de_core} and the DE test titles.
% Adds the Ch11 solution helpers: soln/spart tables, sbs/sbsplain side-by-side, \note, \award, \tagm, \dstart.
% =====================================================================
% de_style.tex: Differential Equations (Ch11) homework / test layout.
% = Ch10 di_style (copied as de_core.tex; header text changed) + the DE
% notation and the framed solution boxes of the v1 DE documents
% (soln / spart, used by the Test worked solutions).
% =====================================================================
\providecommand{\e}{\mathrm{e}}
\providecommand{\cosec}{\operatorname{cosec}}
\providecommand{\dt}{\,\mathrm{d}t}
\providecommand{\abs}[1]{\left|#1\right|}
\providecommand{\dd}[2]{\dfrac{\mathrm{d}#1}{\mathrm{d}#2}}
\providecommand{\dvdx}{\dfrac{\mathrm{d}v}{\mathrm{d}x}}
\providecommand{\dudx}{\dfrac{\mathrm{d}u}{\mathrm{d}x}}
\colorlet{notegrey}{black!60}
\makeatletter
\protected\def\note#1{\ifmmode\text{\small\color{notegrey}\itshape #1}\else{\small\color{notegrey}\itshape #1}\fi}
\newcommand{\award}[1]{\unskip\nobreak\hfill\textbf{\small #1}}
\newcommand{\tagm}[1]{\tag*{\textbf{\small #1}}}
\newcommand{\mfAminus}{\[\begin{aligned}&\int\dfrac{1}{a^{2}-x^{2}}\dx\\&\quad=\dfrac{1}{2a}\ln\abs{\dfrac{a+x}{a-x}}+c\end{aligned}\]}
\newcommand{\mfXminus}{\[\begin{aligned}&\int\dfrac{1}{x^{2}-a^{2}}\dx\\&\quad=\dfrac{1}{2a}\ln\abs{\dfrac{x-a}{x+a}}+c\end{aligned}\]}
\newcommand{\sbs@start}{\if@inlabel\def\sbs@shift{\vspace{-\ht\strutbox}}\else
  \ifvmode\else\par\fi\addvspace{0.4ex}\def\sbs@shift{}\fi\noindent}
\newsavebox{\sbs@box}
\newenvironment{sbs}[1]{%
  \sbs@start\begin{minipage}[t]{#1\linewidth}\vspace{0pt}\sbs@shift\raggedright%
  \def\sbsnext{\end{minipage}\hfill\begin{minipage}[t]{\dimexpr0.97\linewidth-#1\linewidth\relax}\vspace{0pt}\sbs@shift\setlength{\fboxsep}{1.5mm}\begin{lrbox}{\sbs@box}\begin{minipage}[t]{\dimexpr\linewidth-2\fboxsep-1pt\relax}\small\raggedright}%
}{\end{minipage}\end{lrbox}\setlength{\fboxrule}{0.5pt}\fbox{\usebox{\sbs@box}}\end{minipage}\par}
\newenvironment{sbsplain}[1]{%
  \sbs@start\begin{minipage}[t]{#1\linewidth}\vspace{0pt}\sbs@shift\raggedright%
  \def\sbsnext{\end{minipage}\hfill\begin{minipage}[t]{\dimexpr0.97\linewidth-#1\linewidth\relax}\vspace{0pt}\sbs@shift\centering}%
}{\end{minipage}\par}
% ---- framed worked-solution box (v1 DE mark-scheme layout, kept for the Test) ----
\newlength{\solcol}\setlength{\solcol}{15mm}
\newcounter{spartno}
\newcommand{\solskips}{\abovedisplayskip=4pt plus 2pt\belowdisplayskip=4pt plus 2pt\abovedisplayshortskip=2pt\belowdisplayshortskip=3pt}
\newtcolorbox{soln}[1][]{enhanced, breakable, sharp corners,
  colback=white, colframe=black, boxrule=0.6pt,
  boxsep=0pt, left=0pt, right=3mm, top=1.5mm, bottom=2mm,
  before skip=1.2ex, after skip=1.5ex,
  before upper={\setcounter{spartno}{0}\parskip=0.5ex\solskips\@fleqntrue\@mathmargin=0pt\raggedright},
  overlay={\draw[line width=0.4pt] ([xshift=\solcol]frame.north west) -- ([xshift=\solcol]frame.south west);},
  #1}
\newenvironment{spart}[1]{%
  \stepcounter{spartno}%
  \ifnum\value{spartno}>1 \par\vspace{0.2ex}\noindent
    \makebox[\linewidth][l]{\rule{\dimexpr\linewidth+3mm\relax}{0.4pt}}\par\nobreak\vspace{0.3ex}\fi
  \begin{list}{}{\leftmargin=\dimexpr\solcol+2.5mm\relax
    \labelwidth=\solcol \labelsep=2.5mm \itemindent=0pt
    \topsep=0.3ex \parsep=0.5ex \rightmargin=0pt \@beginparpenalty=\@M}%
  \item[{\makebox[\solcol][l]{\strut\hspace{1.5mm}\textbf{#1}}}]%
  \ignorespaces}{\end{list}}
\makeatother
% ---- DE sketches: thin light axes (teacher rule), as the Ch10 'sketch' style ----
\pgfplotsset{dsketch/.style={sketch, axis lines=middle, axis line style={thin, black!45, -{Stealth[length=1.6mm]}},
    ticks=none, clip=false, width=6.5cm, height=4.6cm,
    every axis x label/.style={at={(current axis.right of origin)},anchor=west},
    every axis y label/.style={at={(current axis.above origin)},anchor=south},
    every node/.append style={font=\small}},
  dcurve/.style={black, semithick, smooth, samples=120}}
\tikzset{asym/.style={black, dashed, thin}}
\makeatletter
\newcommand{\dstart}{\leavevmode\vadjust{\nobreak\vskip\dimexpr-\baselineskip-\abovedisplayskip+1pt\relax}}
\newcommand{\sdstart}{\leavevmode\vadjust{\nobreak\vskip\dimexpr-\baselineskip-\abovedisplayskip+4pt\relax}}
\makeatother
